\section{The digit exchange}
Let $\Ag=\{\sum_{j\ge0}\epsilon_j3^j:\epsilon_j\in\{0,1\},\ \epsilon_j=0
\text{ eventually}\}$. If $a+c=2b$ with $a,b,c\in\Ag$, both sides have carry-free ternary expansions. In each digit the only solutions are three zeros or three ones. Thus $a=b=c$. Rankin's set is
\[
 G=\{n\ge1:v_p(n)\in\Ag\text{ for every prime }p\}.
\]
A triple of positive integers is a nonconstant rational-ratio progression precisely when its endpoints $u,w$ and midpoint $v$ satisfy $uw=v^2$ with $u\ne w$. Unique factorization turns this into the simultaneous equalities $v_p(u)+v_p(w)=2v_p(v)$, with at least one nonconstant coordinate. This shows that $G$ is admissible. Write $d(E)=\lim_{X\to\infty}\#(E\cap[1,X])/X$ whenever the limit exists.

\begin{theorem}\label{th:main}
There exists a rational-ratio progression-free $E\subset\N$ with ordinary asymptotic density
\[
                        1.5728\times10^{-727}<d(E)-d(G)<1.5730\times10^{-727}.
\]
For every finite field $\F_q$ there is a subset of nonzero polynomials with ordinary degree-ball density strictly greater than the corresponding greedy density, avoiding every geometric progression with nonconstant ratio in $\F_q(x)$.
\end{theorem}
This interval describes the particular set constructed here. Section~3 first gives a qualitative proof and a transparent one-spoke rule. Section~4 certifies the stated numerical margin. Determining the largest possible density is a separate extremal problem.

\subsection*{A local change with an arbitrary number of primes}
Let $S$ consist of one distinguished prime and $k$ other primes. Number these coordinates $0,1,\ldots,k$. Put $C=\{0,1,3,4\}$, so that $\Ag=C+9\Ag$ uniquely. Call $C^{k+1}$ the old low-digit set. The following table specifies the new low-digit set $V$. Every unlisted coordinate in a displayed pattern has value $1$. In the one-coordinate rows only the distinguished coordinate changes.
\begin{center}
\begin{tabular}{@{}lll@{}}\toprule
Operation & Coordinates shown & Values \\ \midrule
Delete once for each value & $0$ & $1,3,4$ \\
Delete for each $i=1,\ldots,k$ & $(0,i)$ & $(0,4),(1,4),(3,4),(4,4),(4,3)$ \\
Add once & $0$ & $2$ \\
Add for each $i=1,\ldots,k$ & $(0,i)$ & $(2,2)$ \\ \bottomrule
\end{tabular}
\end{center}
The deleted patterns belonged to the old set and are all distinct. Write $x_*$ for the first added vector and $x_i$ for the added vector involving coordinate $i$.

\begin{lemma}\label{lem:digit}
$V$ contains no nonconstant vector triple $u+w=2v$, with unrestricted integral signs in $v-u$.
\end{lemma}
\begin{proof}
Three old vectors satisfy the assertion coordinate by coordinate. Two new vectors in a progression force the third vector's distinguished coordinate to be $2$, so it also is new. All new vectors lie in $\{1,2\}^{k+1}$, where distinct endpoints disagree by one in some coordinate and cannot have an integral midpoint.

It remains to exclude one new vector in a triple with two old vectors. Wherever the new digit equals $1$ and the other two digits belong to $C$, all three must equal $1$, since $C\subset\Ag$. For $x_*$ only the distinguished coordinate remains and the only surviving old pattern with every other digit $1$ has distinguished digit zero. This excludes a triple.

For $x_i$ there are two active coordinates. If $x_i$ is an endpoint, at each active coordinate the ordered old midpoint and other endpoint are either $(1,0)$ or $(3,4)$. If the other endpoint has digit $4$ at $i$, it was deleted, whatever its distinguished digit. If it has digit $0$ at $i$, the midpoint has digit $1$ there and has distinguished digit $1$ or $3$, so the midpoint was deleted. If $x_i$ is the midpoint, the old endpoint digits at $i$ are $0,4$ or $1,3$ in some order. The first possibility loses its digit-$4$ endpoint. In the second possibility the endpoint with digit $1$ at $i$ can survive only with distinguished digit $0$. The opposite endpoint then has distinguished digit $4$ and digit $3$ at $i$, exactly the final deleted pair.
\end{proof}

Define $E(S,V)$ by requiring the exponents on $S$ to have the joint form
\[
       (v_p)_{p\in S}=9h+b,\qquad h\in\Ag^{k+1},\quad b\in V .
\]
At every prime outside $S$ keep the rule $v_p\in\Ag$. If three exponent vectors are in progression, reducing their equation modulo $9$ yields $b_1+b_3=2b_2$ as an \emph{integer} vector equality since both sides lie coordinatewise between $0$ and $8$. Lemma~\ref{lem:digit} forces equality of the three low-digit vectors. Subtraction and division by nine then force equality of the high digits by progression-freeness of $\Ag$. The same conclusion already holds at the outside primes. Unique factorization proves that $E(S,V)$ is progression-free.

\section{Why the density increases}
For a prime with weight $w$ write $H(w)=\sum_{h\in\Ag}w^h$. In the integer problem $w_p=1/p$. The asymptotic probability of a finite exponent pattern $(e_p)_{p\in S}$ is $\prod_{p\in S}(1-w_p)w_p^{e_p}$, by inclusion--exclusion on multiples. Let $c=w_0$, write $w_1,\ldots,w_k$ for the remaining weights, and let $M=c\prod_i w_i$ be the unnormalized weight of the all-ones vector. Relative to $M$, altering one coordinate from $1$ to $a$ multiplies the weight by $w^{a-1}$. The gain in the sum of low-digit weights is $M B$, where
\begin{align}
 B={}&c\left(1+\sum_{i=1}^k w_i\right)-(1+c^2+c^3)\notag\\
    &-c^3\sum_{i=1}^k w_i^2
       -(c^{-1}+1+c^2+c^3)\sum_{i=1}^k w_i^3.                 \label{eq:weight}
\end{align}
The two sums on the second line pay respectively for the deleted pattern $(4,3)$ and the four patterns with digit $4$ at $i$. In the integer example take the distinguished prime to be $2$. Then
\begin{equation}\label{eq:odd}
 B=-\tfrac78+\tfrac12\sum_i p_i^{-1}
             -\tfrac18\sum_i p_i^{-2}-\tfrac{27}{8}\sum_i p_i^{-3}.
\end{equation}
The first sum diverges over the odd primes. Both negative sums converge even over all integers at least two. Therefore a finite choice makes $B>0$. This argument already disproves the conjecture without a numerical search.

\paragraph{Existence of the limit and size of the gain.}
Truncate the old rule at finitely many outside primes and cap the exponents on each selected prime at $L$. The resulting finite-coordinate conditions have densities by counting residue classes. The probability of the discarded high exponents is at most $\sum_{p\in S}p^{-L}$ (and the same argument applies to any finite list under truncation). Ignoring primes $p>Y$ outside $S$ can create an error only at multiples of $p^2$, since $0,1\in\Ag$. Uniformly in $X$, the proportion of such integers up to $X$ is bounded by $\sum_{p>Y}p^{-2}$ when normalized by $X$. Taking the two limits proves existence of the actual density. The resulting difference is exactly
\begin{equation}\label{eq:delta}
 d(E(S,V))-d(G)=M B
 \prod_{p\in S}(1-w_p)H(w_p^9)
 \prod_{p\notin S}(1-w_p)H(w_p).
\end{equation}
The outside product is positive, since each factor is at least $(1-w_p)(1+w_p)=1-w_p^2$. The selected factors are positive as well. Thus the sign of the density difference is the sign of $B$. Since natural density implies the same logarithmic density by partial summation, this construction gives the strict improvement for logarithmic density too.

\paragraph{Polynomial rings and signed rational functions.}
For a monic irreducible $P\in\F_q[x]$ take $w_P=q^{-\deg P}$. Among monic polynomials of degree $n$, the proportion divisible by $a$ is $q^{-\deg a}$ for monic $a$ with $\deg a\le n$, and is zero otherwise. Thus the same finite-dimensional limiting law holds. If $I_q(d)$ counts the degree-$d$ monic irreducibles, then $\sum_{e\mid d}eI_q(e)=q^d$, by factorization of $x^{q^d}-x$. In particular
\[
 I_q(d)\le q^d/d,\qquad
 dI_q(d)\ge q^d-\sum_{j\le d/2,\ j\ge1}q^j\ge q^d/2
 \quad\text{for all sufficiently large }d.
\]
Hence $\sum_P w_P$ diverges and $\sum_P w_P^2$ converges, for every $q$. Fixing a distinguished linear prime now makes \eqref{eq:weight} positive for an appropriate finite list of the remaining irreducibles. The same square tail proves existence of the exact-degree density, and its geometrically weighted average gives the same limit over degree balls. Formula~\eqref{eq:delta} applies with $p$ replaced by $P$.

Include all nonzero scalar multiples of each selected monic polynomial. In a progression with a rational-function ratio, the signed changes in valuations again yield vector midpoint identities. Equal valuations at every monic irreducible force the ratio to lie in $\F_q^\times$, which proves avoidance of every nonconstant rational function, including nonconstant ratios of degree zero. This strengthens the existence request in Question~6.4 of the author manuscript of Asada et al. to ordinary density \cite{asada}. In that question constant ratios are units and are excluded.

\section{Explicit certificates}
The short construction \eqref{eq:odd} already works with all $5{,}132$ odd primes through $50{,}000$. At scale $10^{15}$ the program rounds each summand in the first sum down and in the second and third sums up. It obtains the rigorous normalized bound
\[
 B\ge 543285713107/160000000000000>0.003 .
\]
The selected primorial is below $10^{21599}$ by direct integer comparison. The selected normalizing product is at least $\prod_{2\le n\le 50000}(1-1/n)=1/50000$, the outside product is at least $\prod_{n\ge2}(1-1/n^2)=1/2$, and $H(w^9)\ge1$. Thus this first certificate alone gives $d(E)-d(G)>10^{-21607}$.

In $\F_2[x]$, take $x$ as distinguished prime and all irreducibles of degrees $2$ through $16$ as remaining primes. There are $8{,}799$ selected primes including $x$, with total degree $130{,}485$. The same recurrence for $I_2(d)$ gives
\[
 B=405398067839/70368744177664>1/200 .
\]
Multiplying this normalized gain by the other factors in \eqref{eq:delta} gives the actual density difference.

\subsection*{A larger certified gain using pairs of primes}

For the integer problem a refinement gives the absolute lower bound
\[
                         d(E)-d(G)>10^{-728}.
\]
Use the same base-nine split and pivot prime $2$. Write $D=\{0,1,3,4\}$
for the old low digits. The spokes are the $265$ odd primes through
$1699$. Every new low vector has pivot $2$, with spoke digits $1$ or
$2$. Include the vector with all spoke digits $1$, all vectors with
one spoke digit $2$, and some vectors with two spoke digits $2$. The
selection of pairs is specified below.

For an old low vector, call a spoke active if its digit differs from
$1$. Delete precisely the following old vectors, with $a\in D$ the
pivot digit.
\begin{itemize}
\item With no active spoke, delete the vector if $a\in\{1,3,4\}$.
\item With one active spoke, delete its vector if the spoke digit
is $4$, or if it is $3$ and $a\ne0$.
\item With two active spokes $p<r$, delete their vector if $\{p,r\}$
is selected and their ordered digits belong to
$\{3,4\}\times\{3,4\}$ or equal $(0,4)$. In this line the selection
condition applies to both alternatives.
\end{itemize}
Keep every remaining old vector.

\paragraph{A finite proof of avoidance.}
Two new vectors in a putative progression force the third pivot to be
$2$, and the cube $\{1,2\}^{266}$ contains only constant progressions.
If there is just one new vector, every spoke where it has digit $1$ must
also be $1$ at both old points, since $D$ has no nonconstant arithmetic
progression. It suffices to check the pivot and the $k=0,1,2$ spokes
where the new vector equals $2$. For $k=2$ that pair qualifies for the
two-spoke deletion rule because its new vector was selected. Write $R$
for the test in the three bullets on these ordered coordinates, with
the pair considered selected. For each $v\in D^{k+1}$ form
\[
                  u=4\mathbf{1}-v
             \quad\text{and}\quad
                  u=2v-2\mathbf{1}.
\]
Retain $\{u,v\}$ when $u\in D^{k+1}$. The first choice treats
$2\mathbf{1}$ as the middle, the second as an endpoint with old
middle $v$. The numbers of distinct unordered pairs are respectively
$4,12,40$ for $k=0,1,2$. In every pair $R(u)$ or $R(v)$ holds.
This last finite identity is verified by the loops in
\texttt{check\_spoke\_pairs.py}. At most $4^3$ values of $v$ occur in
one case. This exhausts all signed progressions with one new point.
All-old progressions were already excluded. The base-nine lifting
preserves this conclusion since each digit sum and each doubled digit
is at most eight.

\paragraph{Selecting and pricing each pair.}
Divide the low-digit weights by the weight of the all-ones vector.
Summing over the four possible old pivot digits gives the multiplier
$2+1+1/4+1/8=27/8$. For $p<r$ the net contribution of selecting the
pair is
\[
 J(p,r)=\frac1{2pr}
 -\frac{27}{8}\left[
       (p^{-2}+p^{-3})(r^{-2}+r^{-3})+\frac{p}{r^3}\right].
\]
The last term inside the brackets is the weight of digits $(0,4)$.
Select exactly the pairs with $J(p,r)>0$. Their additions and
deletions have distinct active spoke sets, so their contributions add
without overlaps. The complete normalized gain is
\[
 B=-\frac78+
 \sum_{3\le p\le1699}\left(\frac1{2p}
                   -\frac{11}{8p^2}-\frac{27}{8p^3}\right)
             +\sum_{3\le p<r\le1699}\max\{0,J(p,r)\},
\]
where the indices in both sums are primes. Here $11/8$ is the sum of
the normalized weights of pivot digits $1,3,4$, which are the three
pivot values deleted when the single active spoke has digit $3$.

The integer checker finds $15{,}693$ selected pairs. It bounds $B$ from
below by rounding every summand down to a multiple of $10^{-15}$ and
obtains
\[
                          B\ge\frac{73099512201}{10^{15}}.
\]
For independent reproduction,
$8p^3r^3 J(p,r)=4p^2r^2-27((p+1)(r+1)+p^4)$. Put $P=\prod_{p\le1699}p$. The low-block weight
normalization including the geometric prefactors is $\varphi(P)/P^2$.
The independent higher-digit weights contribute
$\prod_{p\mid P}\sum_{h\in\Ag}p^{-9h}\ge1$. The outside greedy
density is at least $1/\zeta(2)>1/2$. Finally the checker verifies
the integer comparison
\[
       73099512201\,\varphi(P)\,10^{728}>2\cdot10^{15} P^2.
\]
Multiplying these bounds gives the claimed improvement in natural
density. No floating-point arithmetic enters the certificate.

\paragraph{The actual size of the exchange.}
Let $\eta=(d(E)-d(G))/d(G)$. The identity
$H(w)=(1+w)(1+w^3)H(w^9)$ cancels both the outside factors and all high-digit factors when dividing \eqref{eq:delta} by $d(G)$. Thus the \emph{relative} increase is an exact rational number,
\[
 \eta=B\prod_{p\le1699}\frac{p^3}{(p+1)(p^3+1)} .
\]
Computing $B$ from the displayed finite sum as a rational gives
$2.185320021572<10^{727}\eta<2.185320021573$.
We also certify
\begin{equation}\label{eq:numericgain}
              1.5728 <10^{727}\bigl(d(E)-d(G)\bigr)< 1.5730 .
\end{equation}
For the absolute interval, take $T=100000$ and multiply the factors
$(1-p^{-1})(1+p^{-1})(1+p^{-3})(1+p^{-9})$ for $p\le T$.
This retains the first three ternary positions. Relative to this truncated probability, the true $d(G)$ can be lower by at most
$\sum_{n>T}n^{-2}\le1/T$, or higher by at most
$\sum_{n\ge2}n^{-27}\le 2^{-27}+2^{-26}/26$.
The checker bounds the truncated probability by rounding each product outwards to $18$ decimal places. It tests both ends of \eqref{eq:numericgain} by rational comparisons. This completes Theorem~\ref{th:main}.

For an example of the actual rule, let $P$ be the primorial through $1699$. The exchange drops $P$ and admits $2P$, $2pP$ for each spoke, and $2prP$ for each selected pair. The high digits and primes outside the finite list supply all the other integers represented by these patterns. The change in membership therefore has a genuine limiting density even though the numerical gain is exceptionally small.

